\section{Related Work}
\label{section:related}

\subsection{SAX and Related Systems}

Francalaza et.\ al.~\cite{FrancalanzaTP24}, present an alternate formulation of
SAX embedded in adjoint logic. The resulting computational interpretation models
message-passing concurrency: each variable in their system becomes a channel,
whose properties are modeled by its adjoint modality, and each consequent
becomes a provider of said channels~\cite{FrancalanzaTP24}. They present an
implementation of this language in Go~\cite{FrancalanzaTP24}. Logically then,
Pfenning and Pruiksma~\cite{10.1007/978-3-031-35361-1_1} present SAX with
recursion as a unifying framework between message-passing and shared
memory. Specifically, they show that the shared-memory semantics and the
message-passing semantics have a strong bisimulation on well-typed
terms~\cite{10.1007/978-3-031-35361-1_1}.


DeYoung and Pfenning~\cite{DeYoungP23} present SNAX, a modification of the
standard type theory of SAX that accounts for data layout in shared memory
concurrency. It achieves this, in part, by making eligibility a first-class part
of the system, rather than a label, as described above; they also add recursion
to the system~\cite{DeYoungP23}. Ng and Pfenning~\cite{Ng2024MemoryReuse}
present SNAX with reuse, a single-threaded language based on SNAX that models
memory reuse.


Boyland and Pfenning~\cite{BoylandPfenning2025SAX} present an additional
adjoint-logical formulation of SAX with bidirectional typing rules. They further
show a compilation from adjoint natural deduction~\cite{jang2024adjoint} to
their adjoint system, building the groundwork for a compiler of
adjoint-functional programs~\cite{BoylandPfenning2025SAX}. Somayyajula and
Pfenning~\cite{Somayyajula_2023} applies a type refinement system to SAX,
allowing them to reason about the partial correctness of SAX programs. Their
framework also allows for the synthesis of SAX with Hoare
logic~\cite{Somayyajula_2023}.


\subsection{Related Mechanizations}

Pfenning~\cite{Pfenning2000StructuralCutElimination} presents an elegant proof
for cut-elimination in the sequent calculus with the same structure as
in~\cite{DeYoung2020}, and our mechanization. He also presents a formalization
of this proof in
LF~\cite{Pfenning2000StructuralCutElimination}. Crary~\cite{crary2009substructural}
describes a method for encoding linearity in LF without any linear or modal
extensions. We used his techniques to in our representation of SAX in Beluga to
properly represent eligible formulas in our system. These act similarly to
propositions in linear logic as every eligible formula must be used in a SAX
proof. We also drew on the techniques described in~\cite{Sanomech} to properly
represent SAX in Beluga. %%Probably worth adding another sentence or two here, I
%% couldn't fully understand the paper.
